\section{Spectral cutoff mass}
\label{sec:spectral_cutoff_mass}

Let $D$ be a positive semidefinite operator on a finite-dimensional Hilbert
space, and let $P_b=\mathbf 1_{(-\infty,b]}(D)$ be its spectral projection
onto eigenvalues at most $b$. This is the first-moment estimate used in
\cite[Section~7]{OpenAI2026AreaLaw}.
The construction of the replicated defect count and its Hamiltonian energy
bound are separate results.

\begin{theorem}[Operator bound above a spectral threshold]
\label{thm:spectral_cutoff_order}
\lean{Matrix.PosSemidef.smul_one_sub_spectralCutoff_le}
\leanok
For every real $b$,
\begin{align}
 b(I-P_b)&\le D.
 \label{eq:spectral_cutoff_order}
\end{align}
\end{theorem}
\begin{proof}
\leanok
For an eigenvalue $t$ of $D$, positivity gives $t\ge0$.
If $t\le b$, then $b(1-\mathbf1_{t\le b})=0\le t$.
If $t>b$, then this expression equals $b<t$.
Applying the spectral theorem gives the operator inequality.
\end{proof}

\begin{theorem}[Mass below a positive spectral threshold]
\label{thm:spectral_cutoff_mass}
\lean{Matrix.PosSemidef.spectralCutoff_mass_ge}
\leanok
\uses{thm:spectral_cutoff_order}
For $b>0$ and a unit vector $v$,
\begin{align}
 \langle v,P_bv\rangle
 &\ge 1-\frac{\langle v,Dv\rangle}{b}.
 \label{eq:spectral_cutoff_mass}
\end{align}
\end{theorem}
\begin{proof}
\leanok
\uses{thm:spectral_cutoff_order}
Taking the quadratic form of~\eqref{eq:spectral_cutoff_order} at $v$ gives
$b(1-\langle v,P_bv\rangle)\le\langle v,Dv\rangle$.
Division by $b>0$ proves the claim.
\end{proof}
